\centerline{\bf \Huge Параболы и гиперболы}

\problem{1}
График параболы $y=x^2+b x+1$ пересекает ось $O x$ в точках $A$ и $B$, лежащих справа от оси $O y$, и ось $O y$ в точке $C$. Докажите, что описанная окружность треугольника $A B C$ касается оси $O y$.

\problem{2}
(a) \textbf{Критерий вписанности четырехугольника в параболу.} Треугольник $A B C$ вписан в параболу $y=x^2$. Докажите, что точка $D$ лежит на той же параболе тогда и только тогда, когда $k_{A C}-k_{B C}=k_{A D}-k_{B D}$, где через $k_{\ell}$ обозначен угловой коэффициент прямой $\ell$. Для каких еще парабол выполнено такой соотношение?

\noindent
(б) \textbf{Лемма Фусса для парабол.} Две параболы с параллельными осями пересекаются в точках $A$ и $B$. Через точку $A$ проведена прямая, пересекающая первую параболу в точке $C_1$, а вторую - в точке $C_2$. Через точку $B$ прведена прямая, пересекающая первую параболу в точке $D_1$, а вторую - в точке $D_2$. Докажите, что $C_1 D_1 \| C_2 D_2$.

\problem{3}
Две параболы с перпендикулярными осями пересекаются по четырем точкам. Докажите, что эти четыре точки лежат на одной окружности.

\problem{4}
(a) Дана парабола с осью симметрии, параллельной оси $O y$. Через точку $P$ к параболе проведены две секущих. Первая пересекает параболу в точка $A$ и $B$, а вторая пересекает параболу в точках $C$ и $D$. Докажите, что $\mathbf{P A} \cdot \mathbf{P B} =\mathbf{P C} \cdot \mathbf{P D}$, где через $\mathbf{u}$ обозначена длина проекции отрезка $u$ на ось $O x$.

\noindent
(б) Сформулируйте и докажите обратное утверждение.

\textbf{Определение.} Равнобокой гиперболой (или равнобочной, равносторонней, прямоугольной) называется гипербола с перпендикулярными асимптотами.

\problem{5}
Докажите, что если три вершины треугольника лежат на равнобокой гиперболе, то его ортоцентр тоже лежит на этой гиперболе.

\problem{6}
Через вершины треугольника $ABC$ проведена равнобокая гипербола, пересекающая описанную окружность в точке $P$, отличной от $A,$ $B$ и $C$. Докажите, что $P$ и ортоцентр $H$ симметричны относительно центра симметрии гиперболы.


\newpage

\hspace{3mm}

\problem{7}
a) Всякая секущая к равнобокой гиперболе отсекается равными отрезками в пересечении с осями.

\noindent
б) Окружность Эйлера треугольника, вписанного в равнобокую гиперболу, проходит через точку пересечения асимптот

\noindent
в) $\bigl[\textbf{Финал ВСоШ 2024, 10.4}\bigr]$
Дан выпуклый четырёхугольник $A B C D$, в котором $\angle A+\angle D= 90^{\circ}$, его диагонали пересекаются в точке $E$. Прямая $\ell$ пересекает отрезки $A B, C D, A E$ и $E D$ в точках $X, Y, Z$ и $T$ соответственно. Известно, что $A Z=C E$ и $B E=D T$. Докажите, что длина отрезка $X Y$ не больше диаметра описанной окружности треугольника $E T Z$.